\section{System description and EMC-relevant hardware}
The autonomous robot platform consists of two NiMH batteries connected in series and a control system built around a mainboard, four separate H-bridge PCBs, sensors, and wireless communication modules. An adjustable LM2596S-based switching regulator is connected directly after the batteries and provides the supply for the H bridge PCBs. The mainboard contains an STM32-based controller for motor control and sensor acquisition, an ESP32-based controller for higher-level communication, and an LMR36520 buck converter for the 3.3 V logic supply. The platform also includes a DWM1001 UWB module, an LSM9DS1 IMU, four AS5600 wheel-position sensors, and ultrasonic distance sensors, all of which are connected to the mainboard using ribbon cables.



\subsection{Encoders}
The four AS5600 wheel-position sensors are read using 230 Hz PWM outputs. Each sensor is connected to the mainboard by an approximately 20 cm cable carrying 3.3 V, GND, and PWM. For later EMC analysis, the relevant output parameters from the datasheet~[\ref{ref:as5600-datasheet}] are a PWM slew rate of 0.5 to 2 V/$\mu$s and an output current of 0.5 mA.
The corresponding rise-time-based frequency estimate is:
\[
f_r  = \frac{dV/dt}{3.3} = \frac{0.5 \text{ to } 2}{3.3} = 152 \text{ to } 606~\mathrm{kHz}
\]
These values are an indication of the frequency range over which the encoder signal could still have a relevant amplitude.

The wavelength is:
\[
\lambda = \frac{c}{f}
= \frac{299\,792\,458}{230}
= 1\,303\,445.47~\mathrm{m}
\]
We know that the radiated power is
\begin{equation*}
P_{\mathrm{rad}} \propto \left(\frac{L}{\lambda}\right)^2
\tag*{[\ref{ref:paul-emc}]}
\end{equation*}

This means
\[
P_{\mathrm{rad}} \propto \left(\frac{0.2}{1\,303\,445.47}\right)^2 = 2.35 \times 10^{-14}
\]

This is very low, meaning that EMC emissions from the encoder cables are very unlikely.
\subsection{IMU}
The LSM9DS1 IMU is connected to the mainboard through an SPI interface with a ribbon cable that is at maximum 20cm. In the present design, the SPI clock is set by
\[
\frac{170~\mathrm{MHz}}{64} = 2.65625~\mathrm{MHz}
\]

The wavelength is
\[
\lambda = \frac{c}{f}
= \frac{299\,792\,458}{2\,656\,250}
= 112.86~\mathrm{m}
\]


This means
\[
P_{\mathrm{rad}} \propto \left(\frac{0.2}{112.86}\right)^2 = 3.14 \times 10^{-6} = 3.14~\mu\mathrm{W}
\]


The SPI clock frequency is known, but the slew rate of the SPI signals is not currently known. The slew rate must therefore be measured in order to estimate up to which harmonic numbers the SPI interface may still have relevant spectral amplitudes.

\subsection{Ultrasonic}
The HC-SR04 ultrasonic sensor is connected to the mainboard through VCC, GND, Trigger, and Echo. The Trigger input is a TTL pulse of at least 10~$\mu$s, and the Echo output is a TTL pulse whose width is proportional to the measured distance. The module operates from 5~$\mathrm{V}$ with a typical working current of 15~$\mathrm{mA}$. The 40~$\mathrm{kHz}$ ultrasonic burst is generated internally by the module after the Trigger pulse, rather than being directly carried on the interface cable. The datasheet does not specify the rise time, fall time, or slew rate of the Trigger and Echo signals. For later radiated-EMC analysis, these edge rates must therefore be measured on the implemented hardware.




\subsection{Batteries}
The platform is powered by two NiMH batteries connected in series. The batteries themselves are not expected to generate high frequency electromagnetic disturbances, since they are not switching circuits. From an EMC perspective, their main importance is instead that they supply the parts of the system that can generate disturbances, especially the motor drive and the switching regulators.

The battery cables can carry relatively large and rapidly changing currents when the motors change speed or direction. These current changes can create conducted voltage drops and current loops in the power wiring. The battery connection is therefore relevant for conducted EMC, even though the batteries themselves are not an active source of switching noise. Bulk capacitance placed close to the regulator and motor drive inputs can reduce these effects by providing a local energy reservoir, thereby reducing the transient current that must flow through the battery leads. However, bulk capacitance does not eliminate the problem completely, and its effectiveness depends on component choice and physical placement.

The battery wiring can be approximated as a current loop. If the motor current changes
quickly, a voltage disturbance is produced by the parasitic inductance of the wiring:

\[
V_L = L \frac{\mathrm{d}i}{\mathrm{d}t}
\]

For a rough estimate, a straight wire has an inductance in the order of \(1\,\mu\mathrm{H}/\mathrm{m}\). If
the total battery supply and return path is assumed to be \(0.4\,\mathrm{m}\), the loop inductance is
approximately:

\[
L \approx 0.4\,\mu\mathrm{H}
\]

If the motor current changes by \(2\,\mathrm{A}\) in \(1\,\mu\mathrm{s}\), the induced voltage becomes:

\[
V_L = 0.4\,\mu\mathrm{H} \cdot \frac{2\,\mathrm{A}}{1\,\mu\mathrm{s}} = 0.8\,\mathrm{V}
\]

\subsection{Mainboard}

The mainboard is the common PCB for the STM32 controller, ESP32 controller, sensor connectors, UWB connector, and motor-control connections. It also includes the LMR36520 buck converter for the 3.3~V logic supply.
The individual circuits are described in later sections, so the EMC focus here is the routing
on the mainboard and how the signal return currents are handled.

The mainboard carries several different signal types. The AS5600 encoder signals are
captured as PWM signals at approximately \(230\,Hz\). The motor control PWM signals are
approximately \(20\,kHz\). The DWM1001 UWB module communicates over UART at
\(115200\,baud\), and the UART between the ESP32 and STM32 also uses \(115200\,baud\).
The LSM9DS1 IMU is connected over SPI. In the present design the STM32 SPI clock is:

\[
f_{SPI} = \frac{170\,MHz}{64} = 2.65625\,MHz
\]

The nominal frequencies are therefore:

\[
f_{encoder} = 230\,Hz
\]

\[
f_{motorPWM} = 20\,kHz
\]

\[
f_{UART} \approx 115.2\,kHz
\]

\[
f_{SPI} = 2.65625\,MHz
\]

For the mainboard traces themselves, these frequencies are not high enough to make a
short PCB trace an efficient radiator by length alone. This can be checked using the
wavelength relation:

\[
\lambda = \frac{v}{f}
\]

This relation is used when discussing electrical dimensions, where the
important quantity is conductor length compared with wavelength rather than physical
length alone. For the highest nominal mainboard signal frequency,
\(2.65625\,MHz\), the wavelength is:

\[
\lambda_{SPI} =
\frac{3 \cdot 10^8}{2.65625 \cdot 10^6}
= 112.9\,m
\]

For a typical \(50\,mm\) mainboard trace this gives:

\[
\frac{L}{\lambda} =
\frac{0.05}{112.9}
= 4.43 \cdot 10^{-4}
\]

This is far below \(0.1\lambda\), which Paul et al. use as a rule of thumb for when a
structure can be treated as electrically small~[\ref{ref:paul-emc}]. Therefore, the SPI trace
length itself is not expected to be the dominant radiated-emission problem. The encoder
PWM, motor PWM logic signal, and UART traces have even lower nominal frequencies,
so their wavelength-based trace-length risk is lower than the SPI case.

This does not mean that the mainboard layout is irrelevant. Paul et al. also describe that
digital signals contain high-frequency components depending on rise and fall time, and
that PCB design must control return paths and loop areas~[\ref{ref:paul-emc}]. Therefore, the
important mainboard issue is not the nominal frequency alone, but whether the outgoing
trace and its return current stay close together. If a signal trace crosses a weak or broken
ground return, the return current must take a longer path. This increases loop area and
can increase magnetic coupling and radiated emission.

The pre-compliance measurements for the mainboard should therefore focus on signal
quality and coupling, rather than treating the short PCB traces as antennas. 


The main EMC concern on the mainboard is therefore the LMR36520 buck layout rather
than the short logic traces. The buck has a local high-frequency current loop between the
input capacitor, the switching path, ground, and back to the input capacitor. This loop
should be kept small, because loop inductance creates voltage disturbance when the current
changes quickly:

\[
V_L = L \frac{di}{dt}
\]

For this mainboard, the return path is made with copper pours and several vias instead of
a continuous ground plane. This makes via placement important, especially around the
buck. The input capacitor is placed close to the buck, and the buck ground
return to the input capacitor ground through several nearby vias instead of one shared
narrow return path. This reduces common impedance and local ground bounce in the
switching current path.

The SW node copper area should be kept small, because it has the fastest voltage
transitions on the mainboard. The SW node should also be kept away from SPI, UART,
encoder, UWB, and feedback traces. This follows the PCB EMC guidance in Paul et al.,
where return-current path, ground grid/vias, power distribution, and loop area are treated
as key layout factors.
\subsection{H-bridge PCBs}
The robot uses four separate self designed H bridge PCBs for motor driving. Each full
bridge is built from two half bridges, using IRS2008S gate drivers and
IPD220N06L3GATMA1 N channel MOSFETs. The IRS2008S is a high and low side
MOSFET driver, while the IPD220N06L3GATMA1 is a 60 V power MOSFET. This makes
the H bridge PCBs one of the more important EMC risk areas in the AGV, because they
combine fast gate drive, MOSFET switching, motor current, and external motor wiring.

The motor PWM frequency is approximately \(20\,kHz\). This frequency is not the main
radiated concern by itself. The larger EMC risk comes from the switching edges and the
current loop formed by the H bridge, motor cable, motor winding, and return path. These
loops can create radiated emissions, while the same switching current can also create
conducted disturbances on the supply wiring.
\subsection{LM2596S-based switching regulator}
The LM2596S-based switching regulator is a purchased buck module used between the
batteries and the H bridge PCBs and mainboard. Like the LMR36520, the regulator IC
itself is considered a lower-probability EMC problem than the surrounding power wiring
and implementation, since it is a commercial regulator IC rather than a self-designed
switching stage. It is still notable that the module switches at approximately \(150\,kHz\),
which can create conducted ripple on the supply wiring.

\subsection{STM32-based controller}
The STM32-based controller uses an STM32G474RE microcontroller. This is a commercial
microcontroller from STMicroelectronics, not a self-designed digital circuit. According to
the datasheet, the STM32G474RE can operate at up to \(170\,MHz\). This makes it one of
the highest-frequency digital devices in the AGV, and therefore EMC relevant even though
the controller itself is not expected to be the same type of emission source as the motor
drive or switching regulators.

\subsection{ESP32-based controller}
The ESP32 based controller uses an ESP32 WROOM 32 module. The module contains the
ESP32 microcontroller and the RF parts needed for 2.4 GHz WiFi and Bluetooth
communication. In this project, the ESP32 runs at \(240\,MHz\), which makes it one of
the highest frequency digital components in the AGV.

The ESP32 is used for Bluetooth communication and for UART communication with the
STM32 at \(115200\,baud\). The UART interface is not expected to be a dominant emission
source by frequency, but the ESP32 module is still EMC relevant because it contains both
a high frequency digital controller and a 2.4 GHz radio.

The radio regulatory part is outside the scope of this report. In the emission analysis,
Bluetooth activity is therefore treated separately from the unintentional emissions of the
AGV electronics. Emissions around \(2.4\,GHz\) are expected from the ESP32 Bluetooth
radio when it is active, while the motor drive, buck regulators, and wiring are treated as
the relevant sources for unintentional emissions.

\subsection{LMR36520 buck converter}
The LMR36520 is a documented commercial synchronous buck regulator from Texas
Instruments, not a self-designed switching stage. According to the datasheet, it includes
integrated high side and low side MOSFETs, internal compensation, and operates as a
4.2 V to 65 V, 2 A step down converter. The datasheet does not state that the IC itself
meets EU EMC emission limits, and the responsibility for meeting EMC requirements is
still on the final AGV product.

Nevertheless, since the LMR36520 is a commercial IC from a major semiconductor
manufacturer and is intended to be used in many different products, it is reasonable to
assume that the IC itself is less likely to be the main EMC problem than the surrounding
buck implementation. For this project, the more probable EMC problem areas are therefore
the PCB layout around the buck, especially the input loop, SW node, inductor, output
capacitors, and return path.

\subsection{DWM1001 UWB module}
The DWM1001 UWB module is a commercial radio module used for global positioning.
Since the regulatory radio part, such as EU RED compliance for the UWB transmitter, is
outside the scope of this project, the UWB module is not connected during the
pre-compliance emission measurements.